\clearpage\corrigesChapitre{1}{Mesurer et raisonner}
\corrigesSouples
\begin{blocLivre}{cor-1-1}\begin{corrige}{1}{Changer de préfixe}
\begin{questions}
\item $d=\qty{4,50e-4}{m}$ et $\Delta t=\qty{2,5e-3}{s}$. On applique les facteurs $10^{-6}$ et $10^{-3}$.
\item $\qty{7,2e-5}{m}$ : la mantisse est comprise entre 1 et 10.
\end{questions}
\end{corrige}\end{blocLivre}
\begin{blocLivre}{cor-1-2}\begin{corrige}{2}{Un résultat sans unité}
\begin{questions}
\item Le nombre 25 ne précise pas l’unité de volume ; 25 mL et 25 L décrivent des quantités différentes.
\item $V=\qty{25}{mL}=\qty{0,025}{L}$.
\end{questions}
\end{corrige}\end{blocLivre}
\begin{blocLivre}{cor-1-3}\begin{corrige}{3}{Lire un instrument}
\begin{questions}
\item 8,4 cm est une écriture cohérente avec le millimètre. Les zéros supplémentaires annonceraient une précision non justifiée.
\item Le positionnement du zéro et l’angle d’observation peuvent modifier la lecture.
\end{questions}
\end{corrige}\end{blocLivre}
\begin{blocLivre}{cor-1-4}\begin{corrige}{4}{Une moyenne simple}
\begin{questions}
\item $\bar t=(2,1+2,0+2,2+2,1)/4=\qty{2,1}{s}$.
\item Une fluctuation est attendue. Une suppression demande une cause identifiée, par exemple un déclenchement raté, et doit être documentée.
\end{questions}
\end{corrige}\end{blocLivre}
\begin{blocLivre}{cor-1-5}\begin{corrige}{5}{Deux distributions}
\schema[corrige]{histogramme}
\begin{questions}
\item La série A est la plus resserrée autour de 10,0 mL.
\item Non. Des valeurs peuvent être très regroupées autour d’une valeur décalée ; il faut également vérifier l’instrument et comparer à une référence.
\end{questions}
\end{corrige}\end{blocLivre}
\begin{blocLivre}{cor-1-6}\begin{corrige}{6}{Estimer une incertitude}
\begin{questions}
\item $u=0,12/4=\qty{0,03}{mL}$.
\item $\bar V=\qty{24,96}{mL}$ et $u(\bar V)=\qty{0,03}{mL}$. Les deux écritures vont au centième.
\end{questions}
\end{corrige}\end{blocLivre}
\begin{blocLivre}{cor-1-7}\begin{corrige}{7}{Arrondir à la fin}
\begin{questions}
\item La calculatrice donne $v=1,302083\ldots\unit{m.s^{-1}}$.
\item $v=\qty{1,3}{m.s^{-1}}$ ; la durée comporte deux chiffres significatifs.
\end{questions}
\end{corrige}\end{blocLivre}
\begin{blocLivre}{cor-1-8}\begin{corrige}{8}{Comparer à une référence}
\begin{questions}
\item 0,789 est très proche de 0,79 ; 1,00 est éloigné de 0,21, soit 21 fois l’incertitude annoncée.
\item Des espèces ou des mélanges peuvent avoir des masses volumiques proches. Il faut croiser plusieurs indices.
\end{questions}
\end{corrige}\end{blocLivre}
\begin{blocLivre}{cor-1-9}\begin{corrige}{9}{Un quotient trompeur}
\begin{questions}
\item $500\unit{mg.L^{-1}}$ car la masse a été saisie en milligrammes.
\item $\qty{0,50}{g.L^{-1}}$. Convertir l’unité ne change pas la grandeur physique.
\end{questions}
\end{corrige}\end{blocLivre}
\begin{blocLivre}{cor-1-10}\begin{corrige}{10}{Répéter utilement}
\begin{questions}
\item La seconde : elle produit des mesures indépendantes du volume délivré. Recopier une valeur ne crée aucune information nouvelle.
\item Oui. Il faut maîtriser les conditions qui affectent la masse mesurée pour interpréter la dispersion.
\end{questions}
\end{corrige}\end{blocLivre}
\begin{blocLivre}{cor-1-11}\begin{corrige}{11}{Améliorer une méthode}
\begin{questions}
\item Chronométrer plusieurs oscillations puis diviser par leur nombre ; répéter la manipulation. Le délai de réaction pèse moins sur une durée plus longue.
\item $T=8,16/10=\qty{0,816}{s}$ avant choix de l’arrondi selon l’incertitude.
\end{questions}
\end{corrige}\end{blocLivre}
\begin{blocLivre}{cor-1-12}\begin{corrige}{12}{Histogramme et classes}
\demarche{Compter chaque mesure dans une seule classe, puis calculer la moyenne.}
\begin{questions}
\item Les effectifs sont 2, 3 et 1. Chaque valeur est comptée une seule fois.
\item $\bar V=\qty{10,0}{mL}$. L’histogramme montre la répartition et la dispersion que la moyenne seule ne décrit pas.
\end{questions}
\retiens{La moyenne situe la série ; l’histogramme décrit aussi sa dispersion.}
\end{corrige}\end{blocLivre}

